\documentclass[11pt]{article}
\usepackage[notes]{calc}
\calcmodule{8}{6}{Integration Strategies}
\calcxdoc{8.1}
\calcxdoc{8.2}
\calcxdoc{8.3}
\calcxdoc{8.4}
\calcxdoc{8.5}

\begin{document}

\section{Integration Strategies}

\begin{guide}
Sections 8.1 to 8.5 taught one method at a time, so the method was always announced by the title of the section. A real integral does not come with a title. This section is about the question that comes first: \emph{which method should I try?} Differentiation is mechanical, because the rules of Calculus~I always apply. Integration is a craft: the same integral can often be done in several ways, one of them is usually much shorter, and some integrals cannot be done in closed form at all. We give a decision procedure (Procedure~\ref{proc:strategy}), a table of patterns in the integrand, and a collection of examples that need a rewrite, a dead end and a fresh start, or two methods in a row. We prove the facts that make the craft safe: how to \emph{check} an answer, why two correct answers may look different (Theorem~\ref{thm:check}), why a substitution can turn a root or an exponential into a rational function (Theorem~\ref{thm:rational}), and how to \emph{recognize} an integral with no elementary antiderivative (Theorems~\ref{thm:liouville} and~\ref{thm:norat}).
\end{guide}

\textbf{By the end of this section you should be able to:}
\begin{itemize}
  \item apply the decision procedure: rewrite, test for a substitution, recognize the pattern, try the first steps of more than one method, and abandon a method that makes the integral worse;
  \item combine two or three methods in sequence (a substitution followed by parts, parts followed by partial fractions, a substitution followed by a trigonometric substitution);
  \item use a rationalizing substitution for $(ax+b)^{1/n}$ and for $e^x$;
  \item check every antiderivative by differentiating it, and say why a table or a computer algebra result is a claim and not a proof;
  \item recognize integrals such as $\int e^{-x^2}\,dx$ that have no elementary antiderivative, say what that statement means and what it does not mean, and give a numerical value or a bound for the corresponding definite integral.
\end{itemize}

% ---------------------------------------------------------------------
\subsection*{A decision procedure}

An integral is attacked in a fixed order of questions. The order is not a law of nature; it is a list of the cheapest things to try first.

\begin{proc}{Choosing a method}{strategy}
\begin{enumerate}[nosep]
  \item \emph{Is it a basic formula?} Compare with the table of Section~8.1 (formulas (1)--(17)). If a constant factor or a linear argument $ax+b$ is the only difference, you are done.
  \item \emph{Can it be rewritten?} Split a fraction, divide polynomials, use an identity, complete the square, multiply by a conjugate or by $1$ (Procedure~\ref{8.1-proc:rewrite}). Rewriting costs nothing and often makes step 3 or 4 obvious.
  \item \emph{Is there a substitution?} Look for a function $g(x)$ such that $g'(x)$, up to a constant factor, is also in the integrand: the inside of a power, a root, an exponential or a trigonometric function; a denominator whose derivative is the numerator (Procedure~\ref{8.1-proc:subst}).
  \item \emph{Recognize the pattern} (the table below) and use the method of the matching section: parts (8.2), trigonometric integrals (8.3), trigonometric substitutions (8.4), partial fractions (8.5).
  \item \emph{Try the first steps of more than one method.} Stop as soon as the new integral is \emph{no simpler}: higher powers, more factors, a return to the start. A dead end is information, not failure.
  \item \emph{Combine.} After a substitution, parts or a division the new integrand is a new problem: start again at question 1.
  \item \emph{If every route fails, ask whether an elementary antiderivative exists at all} (Theorems~\ref{thm:liouville} and~\ref{thm:norat}); a definite integral can still be computed numerically.
  \item \emph{Check.} Differentiate an indefinite answer; compare a definite answer with a numerical value or a rough estimate of sign and size.
\end{enumerate}
\end{proc}

\begin{center}
\renewcommand{\arraystretch}{1.3}
\begin{tabular}{p{2.3in}p{2.75in}l}
\textbf{Pattern in the integrand} & \textbf{First method} & \textbf{Section}\\ \hline
$f(g(x))\,g'(x)$ & substitution $u=g(x)$ & 8.1\\
polynomial $\times$ $e^{ax}$, $\sin ax$, $\cos ax$ & parts, $u=$ the polynomial (repeat, or tabular) & 8.2\\
$\ln x$, $\tan^{-1}x$, $\sin^{-1}x$, polynomial $\times$ one of them & parts, $dv=$ the polynomial part & 8.2\\
$e^{ax}\sin bx$, $e^{ax}\cos bx$ & parts twice and solve (cyclic) & 8.2\\
$\sin^mx\cos^nx$ & odd power: save one factor, substitute; even powers: identities & 8.3\\
$\tan^mx\sec^nx$, $\cot^mx\csc^nx$ & $u=\tan x$ or $u=\sec x$ (or reduction formula) & 8.3\\
$\sqrt{a^2-x^2}$, $\sqrt{a^2+x^2}$, $\sqrt{x^2-a^2}$ & trigonometric substitution, \emph{unless} a plain substitution works & 8.4\\
$ax^2+bx+c$ in a root or a denominator & complete the square & 8.1, 8.4\\
rational function & divide if improper; substitution if the numerator is the derivative; else partial fractions & 8.5\\
$(ax+b)^{1/n}$, or $x$ only inside $e^x$ & rationalizing substitution (Theorem~\ref{thm:rational}) & this section\\
\end{tabular}
\end{center}

% ---------------------------------------------------------------------
\subsection*{Checking an answer, and why two answers can differ}

Every method in this chapter produces a \emph{formula}. Whether the formula is correct can be settled in one line, independently of how it was found.

\begin{thm}{Checking and comparing antiderivatives}{check}
Let $f$ be defined on an interval $I$.
\begin{enumerate}[label=(\arabic*)]
  \item A differentiable function $F$ is an antiderivative of $f$ on $I$ if and only if $F'(x)=f(x)$ for every $x$ in $I$. So any proposed answer, from any source, can be tested by differentiating it.
  \item If $F$ and $G$ are antiderivatives of $f$ on $I$, then $F-G$ is constant on $I$. Two correct answers obtained by different methods can therefore look different, but they differ by a constant.
\end{enumerate}
\end{thm}
\begin{proof}
(1) This is the definition of an antiderivative. (2) Put $H=F-G$. Then $H'=F'-G'=f-f=0$ on $I$. Let $p<q$ be two points of $I$. By the Mean Value Theorem there is a number $s$ between $p$ and $q$ with $H(q)-H(p)=H'(s)(q-p)=0$. So $H(q)=H(p)$ for all $p,q$ in $I$, and $H$ is constant.
\end{proof}

Part (2) needs an \emph{interval}. For $\frac1x$, the function that equals $\ln x+5$ for $x>0$ and $\ln(-x)-3$ for $x<0$ is an antiderivative on $x\ne0$, yet it is not $\ln\lvert x\rvert$ plus one constant: the constant may change from one interval to the next. The usual ``$+C$'' means one constant for each interval on which the integrand is defined.

\begin{warn}[A table or a computer result is not a proof]
A table of integrals, a computer algebra system or a classmate gives you a \emph{claim}: ``$F$ is an antiderivative of $f$''. It becomes a proven fact only when you differentiate $F$ and get $f$ (Theorem~\ref{thm:check}(1)). Tables and programs may leave out absolute values, give a formula valid only on part of the domain, or write the answer in a form that differs from yours by a constant (Theorem~\ref{thm:check}(2)). In this course every integral in the examples is checked by differentiation, and a table or computer result is used only as a way to guess an $F$ whose derivative we then compute.
\end{warn}

% ---------------------------------------------------------------------
\subsection*{Rewrite, substitute, recognize}

\begin{exmp}{Rewrite the integrand}{m1}
Evaluate $\displaystyle\int\frac{1+\cos x}{\sin^2x}\,dx$.

\textbf{Solution.} Split the fraction and use $\frac1{\sin x}=\csc x$ and $\frac{\cos x}{\sin x}=\cot x$:
\[\frac{1+\cos x}{\sin^2x}=\frac1{\sin^2x}+\frac1{\sin x}\cdot\frac{\cos x}{\sin x}=\csc^2x+\csc x\cot x .\]
Both terms are in the table of Section~8.1 (formulas (6) and (8) with $a=1$), so on each interval on which $\sin x\ne0$
\[\int\frac{1+\cos x}{\sin^2x}\,dx=-\cot x-\csc x+C .\]
%@ m1a | ANT(-cot(x)-csc(x),(1+cos(x))/sin(x)**2,[Q(1,2),Q(1,1),Q(5,2),-Q(7,2)]) == 0 | =-\cot x-\csc x+C .
\emph{Check.} $\frac d{dx}(-\cot x-\csc x)=\csc^2x+\csc x\cot x$, which is the integrand by the split above.

\emph{Another form.} By the half-angle identity $\cot\frac x2=\frac{1+\cos x}{\sin x}=\csc x+\cot x$, the answer is also $-\cot\frac x2+C$. The two forms are the same function, so by Theorem~\ref{thm:check}(2) the constant is $0$.
%@ m1b | sum(NUM(-cot(v)-csc(v),-cot(v/2)) for v in (Q(1,2),Q(1,1),Q(5,2))) == 0 | the constant is $0$.

\emph{A route that goes in a circle.} If instead we write $\sin^2x=(1-\cos x)(1+\cos x)$ and cancel, the integral becomes $\int\frac{dx}{1-\cos x}$. This is correct, but it is not in our table. Multiplying top and bottom by $1+\cos x$ to get rid of the denominator gives $\frac{1+\cos x}{1-\cos^2x}=\frac{1+\cos x}{\sin^2x}$, the integral we started with. When a rewrite returns to the start, discard it and choose another.
\end{exmp}

\begin{exmp}{A substitution after a rewrite}{m2}
Evaluate $\displaystyle\int\frac{x^3}{\sqrt{25-x^8}}\,dx$ on $\lvert x\rvert<\sqrt[4]5$.

\textbf{Solution.} The root contains $25-x^8=25-(x^4)^2$, and the derivative of $x^4$ is $4x^3$, which is the numerator up to the factor $4$. So rewrite the denominator and put $u=x^4$, $du=4x^3\,dx$, $x^3\,dx=\frac14\,du$:
\[\int\frac{x^3\,dx}{\sqrt{25-(x^4)^2}}=\frac14\int\frac{du}{\sqrt{25-u^2}}=\frac14\sin^{-1}\frac u5+C=\frac14\sin^{-1}\frac{x^4}5+C,\]
using formula (12) of Section~8.1 with $a=5$. (A trigonometric substitution $x^4=5\sin\theta$ would work too, but the plain substitution does the job in one step, as the table says.)
%@ m2 | ANT(asin(x**4/5)/4,x**3/sqrt(25-x**8),[Q(1,2),Q(1,1),Q(3,2),-Q(1,2)]) == 0 | =\frac14\sin^{-1}\frac u5+C=\frac14\sin^{-1}\frac{x^4}5+C,\]
\emph{Check.} $\frac d{dx}\frac14\sin^{-1}\frac{x^4}5=\frac14\cdot\frac1{\sqrt{1-x^8/25}}\cdot\frac{4x^3}5=\frac{x^3}{5\sqrt{(25-x^8)/25}}=\frac{x^3}{\sqrt{25-x^8}}$, because $5\sqrt{\frac{25-x^8}{25}}=\sqrt{25-x^8}$.
\end{exmp}

\begin{exmp}{Identify the method, then finish}{m3}
For each integral decide on a method from the pattern in the integrand, and evaluate it.

(a) $\displaystyle\int\frac{5-3x^2}{x^3-5x}\,dx$. \quad (b) $\displaystyle\int\ln\left(4+x^2\right)dx$. \quad (c) $\displaystyle\int\sin^3x\cos^2x\,dx$.

\textbf{Solution.} (a) \emph{A rational function, but test the substitution first.} The derivative of the denominator is $3x^2-5$, and the numerator is $5-3x^2=-(3x^2-5)$. Let $u=x^3-5x$, $du=(3x^2-5)\,dx$, so $(5-3x^2)\,dx=-du$:
\[\int\frac{5-3x^2}{x^3-5x}\,dx=-\int\frac{du}u=-\ln\lvert u\rvert+C=-\ln\left\lvert x^3-5x\right\rvert+C .\]
%@ m3a | ANT(-log(sqrt((x**3-5*x)**2)),(5-3*x**2)/(x**3-5*x),[Q(1,2),Q(2,1),Q(3,1),-Q(1,1)]) == 0 | =-\ln\left\lvert x^3-5x\right\rvert+C .\]
Partial fractions (Section~8.5) would give the same function after factoring $x^3-5x=x(x-\sqrt5)(x+\sqrt5)$: the three covered-up numerators are $A=\frac5{-5}=-1$ at $x=0$, and $-1$ at each of $x=\pm\sqrt5$ (check: $\frac{5-15}{\sqrt5\cdot2\sqrt5}=-1$), and $-\ln\lvert x\rvert-\ln\lvert x-\sqrt5\rvert-\ln\lvert x+\sqrt5\rvert=-\ln\lvert x^3-5x\rvert$. Same answer, far more work.
%@ m3a2 | simplify(apart((5-3*x**2)/(x**3-5*x),x)-(-1/x-1/(x-sqrt(5))-1/(x+sqrt(5)))) == 0 | Same answer, far more work.

(b) \emph{A single transcendental function with no polynomial factor:} parts with $dv=dx$ (Section~8.2, as for $\ln x$). Let $u=\ln(4+x^2)$, $dv=dx$, so $du=\frac{2x}{4+x^2}\,dx$ and $v=x$:
\[\int\ln\left(4+x^2\right)dx=x\ln\left(4+x^2\right)-\int\frac{2x^2}{4+x^2}\,dx .\]
The new integrand is improper (equal degrees), so divide: $2x^2=2(x^2+4)-8$, hence $\frac{2x^2}{x^2+4}=2-\frac8{x^2+4}$, and $\int\frac8{x^2+4}\,dx=8\cdot\frac12\tan^{-1}\frac x2=4\tan^{-1}\frac x2$ (formula (11), $a=2$). Therefore
\[\int\ln\left(4+x^2\right)dx=x\ln\left(4+x^2\right)-2x+4\tan^{-1}\frac x2+C .\]
%@ m3b | ANT(x*log(4+x**2)-2*x+4*atan(x/2),log(4+x**2),[Q(1,2),Q(3,2),Q(5,2),-Q(2,1)]) == 0 | =x\ln\left(4+x^2\right)-2x+4\tan^{-1}\frac x2+C .\]

(c) \emph{Powers of sine and cosine, odd power of sine} (Section~8.3): save one factor $\sin x$ and write $\sin^2x=1-\cos^2x$. Let $u=\cos x$, $du=-\sin x\,dx$:
\begin{align*}\int\sin^3x\cos^2x\,dx&=\int(1-\cos^2x)\cos^2x\,\sin x\,dx=-\int(1-u^2)u^2\,du\\&=-\frac{u^3}3+\frac{u^5}5+C=-\frac{\cos^3x}3+\frac{\cos^5x}5+C .\end{align*}
%@ m3c | ANT(-cos(x)**3/3+cos(x)**5/5,sin(x)**3*cos(x)**2,[Q(1,2),Q(3,2),Q(5,2)]) == 0 | =-\frac{\cos^3x}3+\frac{\cos^5x}5+C .\end{align*}
\emph{Check of (b) and (c).} $\frac d{dx}\left(x\ln(4+x^2)-2x+4\tan^{-1}\frac x2\right)=\ln(4+x^2)+\frac{2x^2}{4+x^2}-2+\frac{4\cdot\frac12}{1+x^2/4}=\ln(4+x^2)+\frac{2x^2}{4+x^2}-2+\frac8{4+x^2}=\ln(4+x^2)$, because $\frac{2x^2+8}{4+x^2}=2$. And $\frac d{dx}\left(-\frac{\cos^3x}3+\frac{\cos^5x}5\right)=\cos^2x\sin x-\cos^4x\sin x=\sin x\cos^2x(1-\cos^2x)=\sin^3x\cos^2x$.
\end{exmp}

\begin{kc}{Choose the first move}{kc1}
Name the first technique you would try for each integral (do not evaluate): (a) $\displaystyle\int x^2e^{x^3}\,dx$; (b) $\displaystyle\int xe^{3x}\,dx$; (c) $\displaystyle\int\frac{x+1}{x^2+2x+5}\,dx$; (d) $\displaystyle\int\frac{dx}{x^2\sqrt{x^2+9}}$.
\end{kc}

\begin{exmp}{A dead end, and a fresh start}{m4}
Evaluate $\displaystyle\int x^5e^{x^3}\,dx$.

\textbf{Solution.} \emph{First attempt: parts.} The integrand is a power of $x$ times an exponential, a pattern that usually calls for parts. We cannot integrate $e^{x^3}$, so take $u=e^{x^3}$ and $dv=x^5\,dx$: $du=3x^2e^{x^3}dx$, $v=\frac{x^6}6$. Then
\[\int x^5e^{x^3}\,dx=\frac{x^6e^{x^3}}6-\frac12\int x^8e^{x^3}\,dx .\]
%@ m4a | simplify(diff(x**6*exp(x**3)/6,x)-(x**5*exp(x**3)+x**8*exp(x**3)/2)) == 0 | =\frac{x^6e^{x^3}}6-\frac12\int x^8e^{x^3}\,dx .\]
The new integral has the power $x^8$ in place of $x^5$: it is \emph{worse}, so we stop and return to the start.

\emph{Second attempt: substitution.} The exponent $x^3$ has derivative $3x^2$, and $x^5=x^3\cdot x^2$. Let $u=x^3$, $du=3x^2\,dx$, so $x^2\,dx=\frac13\,du$:
\[\int x^5e^{x^3}\,dx=\int x^3\,e^{x^3}\,x^2\,dx=\frac13\int u\,e^u\,du .\]
Now parts is the right tool (Section~8.2): $u=u$, $dv=e^u\,du$ gives $\int ue^u\,du=ue^u-\int e^u\,du=(u-1)e^u$. Therefore
\[\int x^5e^{x^3}\,dx=\frac13\left(x^3-1\right)e^{x^3}+C .\]
%@ m4 | ANT((x**3-1)*exp(x**3)/3,x**5*exp(x**3),[Q(1,2),Q(1,1),-Q(3,2),Q(2,1)]) == 0 | =\frac13\left(x^3-1\right)e^{x^3}+C .\]
\emph{Check.} $\frac d{dx}\left[\frac13(x^3-1)e^{x^3}\right]=\frac13\left[3x^2e^{x^3}+(x^3-1)\cdot3x^2e^{x^3}\right]=x^2e^{x^3}\left(1+x^3-1\right)=x^5e^{x^3}$.
\end{exmp}

\begin{think}{Was the first attempt wrong?}{deadend}
In Example~\ref{ex:m4} the first attempt (parts) was not an error: the integral $\int x^8e^{x^3}\,dx$ can also be done with $u=x^3$. Why was it still the wrong start, and what could you look at, \emph{before} integrating, to see this?
\end{think}

\begin{exmp}{Three methods in a row}{m5}
Find the area under $y=e^x\sqrt{4+e^{2x}}$ for $0\le x\le\ln3$, that is, evaluate $\displaystyle\int_0^{\ln3}e^x\sqrt{4+e^{2x}}\,dx$.

\textbf{Solution.} \emph{Substitution.} The factor $e^x$ is the derivative of the inside function $e^x$, and $e^{2x}=(e^x)^2$. Let $u=e^x$, $du=e^x\,dx$. The limits change: $x=0$ gives $u=1$ and $x=\ln3$ gives $u=3$. So the integral is $\int_1^3\sqrt{4+u^2}\,du$ (the second picture of Figure~1).

\emph{Trigonometric substitution} (Section~8.4). The root $\sqrt{4+u^2}$ has the form $\sqrt{a^2+u^2}$ with $a=2$. Let $u=2\tan\theta$, $-\frac\pi2<\theta<\frac\pi2$, so $du=2\sec^2\theta\,d\theta$ and $\sqrt{4+u^2}=\sqrt{4\sec^2\theta}=2\sec\theta$ (because $\sec\theta>0$ on this interval). Then
\[\int\sqrt{4+u^2}\,du=\int2\sec\theta\cdot2\sec^2\theta\,d\theta=4\int\sec^3\theta\,d\theta .\]

\emph{A reduction formula} (Section~8.3). The secant reduction formula with $n=3$ gives $\int\sec^3\theta\,d\theta=\frac12\sec\theta\tan\theta+\frac12\int\sec\theta\,d\theta$, and $\int\sec\theta\,d\theta=\ln\lvert\sec\theta+\tan\theta\rvert$ (Theorem~\ref{8.1-thm:trig4}). So
\[4\int\sec^3\theta\,d\theta=2\sec\theta\tan\theta+2\ln\lvert\sec\theta+\tan\theta\rvert+C .\]

\emph{Back to $u$} (Theorem~\ref{8.4-thm:back}): $\tan\theta=\frac u2$ and $\sec\theta=\frac{\sqrt{4+u^2}}2>0$. Then $2\sec\theta\tan\theta=\frac{u\sqrt{4+u^2}}2$ and $2\ln(\sec\theta+\tan\theta)=2\ln\frac{u+\sqrt{4+u^2}}2=2\ln\left(u+\sqrt{4+u^2}\right)-2\ln2$; the constant $-2\ln2$ is absorbed into $C$. So an antiderivative is
\[F(u)=\frac12\,u\sqrt{4+u^2}+2\ln\left(u+\sqrt{4+u^2}\right).\]
%@ m5a | ANT(x*sqrt(4+x**2)/2+2*log(x+sqrt(4+x**2)),sqrt(4+x**2),[Q(0,1),Q(1,1),Q(3,1),Q(5,2)]) == 0 | F(u)=\frac12\,u\sqrt{4+u^2}+2\ln\left(u+\sqrt{4+u^2}\right).\]
\emph{Check.} $\frac{d}{du}\ln\left(u+\sqrt{4+u^2}\right)=\frac{1+u/\sqrt{4+u^2}}{u+\sqrt{4+u^2}}=\frac{1}{\sqrt{4+u^2}}$, so $F'(u)=\frac12\sqrt{4+u^2}+\frac{u^2}{2\sqrt{4+u^2}}+\frac2{\sqrt{4+u^2}}=\frac{(4+u^2)+u^2+4}{2\sqrt{4+u^2}}=\frac{4+u^2}{\sqrt{4+u^2}}=\sqrt{4+u^2}$.

\emph{Evaluate.} $F(3)=\frac{3\sqrt{13}}2+2\ln\left(3+\sqrt{13}\right)$ and $F(1)=\frac{\sqrt5}2+2\ln\left(1+\sqrt5\right)$, so
\[\text{Area}=F(3)-F(1)=\frac{3\sqrt{13}-\sqrt5}2+2\ln\frac{3+\sqrt{13}}{1+\sqrt5}\approx5.7174 .\]
%@ m5b | abs(((3*sqrt(13)-sqrt(5))/2+2*log((3+sqrt(13))/(1+sqrt(5)))).evalf()-5.7174) < 1e-4 and (lambda F: NUM(F(3)-F(1),(3*sqrt(13)-sqrt(5))/2+2*log((3+sqrt(13))/(1+sqrt(5)))))(lambda u: u*sqrt(4+u**2)/2+2*log(u+sqrt(4+u**2))) == 0 | \text{Area}=F(3)-F(1)=\frac{3\sqrt{13}-\sqrt5}2+2\ln\frac{3+\sqrt{13}}{1+\sqrt5}\approx5.7174 .\]
\emph{Check by a numerical value.} A numerical evaluation of the original integral gives $5.71740$, in agreement. The picture shows why the answer is believable: the region has height from $\sqrt5\approx2.24$ to $3\sqrt{13}\approx10.8$ over a base of length $\ln3\approx1.10$, so the area lies between $2.24\cdot1.10\approx2.5$ and $10.8\cdot1.10\approx11.9$.
%@ m5c | NUM(Integral(exp(x)*sqrt(4+exp(2*x)),(x,0,log(3))),(3*sqrt(13)-sqrt(5))/2+2*log((3+sqrt(13))/(1+sqrt(5)))) == 0 | so the area lies between $2.24\cdot1.10\approx2.5$ and $10.8\cdot1.10\approx11.9$.

\begin{center}
\begin{tikzpicture}
  \begin{scope}[x=2.4cm,y=0.3cm]
    \draw[->] (-0.1,0) -- (1.35,0) node[right] {$x$};
    \draw[->] (0,-0.4) -- (0,12.5) node[above] {$y$};
    \fill[blue!15] (0,0) -- plot[domain=0:1.0986,samples=50,variable=\t] ({\t},{exp(\t)*sqrt(4+exp(2*\t))}) -- (1.0986,0) -- cycle;
    \draw[very thick,blue!70!black,domain=0:1.0986,samples=60,variable=\t] plot ({\t},{exp(\t)*sqrt(4+exp(2*\t))});
    \fill (0,2.236) circle (1.5pt) node[left,font=\small] {$(0,\sqrt5)$};
%@ f1-a | NUM(exp(0)*sqrt(4+exp(0)),sqrt(5)) == 0 | \fill (0,2.236) circle (1.5pt) node[left,font=\small] {$(0,\sqrt5)$};
    \fill (1.0986,10.8167) circle (1.5pt) node[above left,font=\small] {$(\ln3,\,3\sqrt{13})$};
%@ f1-b | NUM(exp(log(3))*sqrt(4+exp(2*log(3))),3*sqrt(13)) == 0 and abs(float(3*sqrt(13))-10.8167) < 1e-4 | \fill (1.0986,10.8167) circle (1.5pt) node[above left,font=\small] {$(\ln3,\,3\sqrt{13})$};
    \node[below,font=\small] at (1.0986,0) {$\ln3$};
    \node[font=\small] at (0.55,4) {area $\approx5.717$};
    \node[font=\small\bfseries] at (0.6,-4.2) {(a) in $x$};
  \end{scope}
  \begin{scope}[xshift=6.3cm,x=1.0cm,y=0.9cm]
    \draw[->] (-0.1,0) -- (3.6,0) node[right] {$u$};
    \draw[->] (0,-0.1) -- (0,4.2) node[above] {$y$};
    \fill[blue!15] (1,0) -- plot[domain=1:3,samples=50,variable=\t] ({\t},{sqrt(4+\t*\t)}) -- (3,0) -- cycle;
    \draw[very thick,blue!70!black,domain=0:3.4,samples=60,variable=\t] plot ({\t},{sqrt(4+\t*\t)});
    \fill (1,2.236) circle (1.5pt) node[above left,font=\small] {$(1,\sqrt5)$};
%@ f1-c | NUM(sqrt(4+1**2),sqrt(5)) == 0 and exp(0) == 1 | \fill (1,2.236) circle (1.5pt) node[above left,font=\small] {$(1,\sqrt5)$};
    \fill (3,3.606) circle (1.5pt) node[above left,font=\small] {$(3,\sqrt{13})$};
%@ f1-d | NUM(sqrt(4+3**2),sqrt(13)) == 0 and exp(log(3)) == 3 | \fill (3,3.606) circle (1.5pt) node[above left,font=\small] {$(3,\sqrt{13})$};
    \node[below,font=\small] at (1,0) {$u=1$};
    \node[below,font=\small] at (3,0) {$u=3$};
    \node[font=\small] at (2,1.2) {area $\approx5.717$};
    \node[font=\small\bfseries] at (1.8,-1.3) {(b) in $u=e^x$};
  \end{scope}
\end{tikzpicture}
\end{center}
\begin{center}\small Figure 1. (a) The region under $y=e^x\sqrt{4+e^{2x}}$ for $0\le x\le\ln3$. (b) After the substitution $u=e^x$ the same area is the area under $y=\sqrt{4+u^2}$ for $1\le u\le3$. The shapes differ (note the different vertical scales: $0.3$ cm per unit in (a), $0.9$ cm per unit in (b)), but both regions have area $\frac{3\sqrt{13}-\sqrt5}2+2\ln\frac{3+\sqrt{13}}{1+\sqrt5}\approx5.717$.\end{center}
\end{exmp}

\begin{exmp}{Parts, then division, then partial fractions}{m6}
Evaluate $\displaystyle\int\ln\left(x^2+x-6\right)dx$ on $x>2$ (the integrand is defined where $x^2+x-6=(x-2)(x+3)>0$, that is, for $x>2$ and for $x<-3$).

\textbf{Solution.} \emph{Parts} (as for $\ln x$): $u=\ln(x^2+x-6)$, $dv=dx$, so $du=\frac{2x+1}{x^2+x-6}\,dx$ and $v=x$:
\[\int\ln\left(x^2+x-6\right)dx=x\ln\left(x^2+x-6\right)-\int\frac{x(2x+1)}{x^2+x-6}\,dx .\]
\emph{Division} (the numerator $2x^2+x$ has the same degree as the denominator): $2x^2+x=2(x^2+x-6)+(12-x)$, so $\frac{2x^2+x}{x^2+x-6}=2+\frac{12-x}{(x-2)(x+3)}$.

\emph{Partial fractions} (Section~8.5): $\frac{12-x}{(x-2)(x+3)}=\frac A{x-2}+\frac B{x+3}$ with $A=\frac{12-2}{2+3}=2$ (cover up $x-2$ at $x=2$) and $B=\frac{12+3}{-3-2}=-3$ (at $x=-3$). Check at $x=0$: $\frac{12}{-6}=-2$ and $\frac2{-2}+\frac{-3}3=-1-1=-2$. So
\[\int\frac{x(2x+1)}{x^2+x-6}\,dx=\int\left(2+\frac2{x-2}-\frac3{x+3}\right)dx=2x+2\ln\lvert x-2\rvert-3\ln\lvert x+3\rvert,\]
and therefore
\[\int\ln\left(x^2+x-6\right)dx=x\ln\left(x^2+x-6\right)-2x-2\ln(x-2)+3\ln(x+3)+C\qquad(x>2).\]
%@ m6 | ANT(x*log(x**2+x-6)-2*x-2*log(sqrt((x-2)**2))+3*log(sqrt((x+3)**2)),log(x**2+x-6),[Q(3,1),Q(5,2),Q(7,1),-Q(4,1),-Q(5,1)]) == 0 | -2x-2\ln(x-2)+3\ln(x+3)+C\qquad(x>2).\]
\emph{Check.} The derivative is $\ln(x^2+x-6)+\frac{x(2x+1)}{x^2+x-6}-2-\frac2{x-2}+\frac3{x+3}$, and $\frac{x(2x+1)}{x^2+x-6}-2=\frac{12-x}{(x-2)(x+3)}=\frac2{x-2}-\frac3{x+3}$ cancels the last two terms. (For $x<-3$ the same formula holds with $\ln\lvert x-2\rvert$ and $\ln\lvert x+3\rvert$.)
\end{exmp}

% ---------------------------------------------------------------------
\subsection*{Making an integrand rational}

A substitution that is not suggested by an inside function can still be the right idea when it removes a root or an exponential. After the substitution the integrand is a rational function, and Section~8.5 finishes the job.

\begin{thm}{Rationalizing substitutions}{rational}
Let $R(p,q)$ be a rational function of two variables.
\begin{enumerate}[label=(\arabic*)]
  \item Let $n\ge2$ be an integer, $a\ne0$, and let $I$ be an open interval on which $ax+b>0$ (and on which the integrand is continuous). With $u=(ax+b)^{1/n}$,
  \[\int R\left(x,(ax+b)^{1/n}\right)dx=\left[\int R\left(\frac{u^n-b}a,\,u\right)\frac na\,u^{n-1}\,du\right]_{u=(ax+b)^{1/n}},\]
  and the integrand on the right is a rational function of $u$.
  \item With $u=e^x$, $\displaystyle\int R(e^x)\,dx=\left[\int\frac{R(u)}u\,du\right]_{u=e^x}$, and the integrand on the right is a rational function of $u$.
\end{enumerate}
\end{thm}
\begin{proof}
We use Theorem~\ref{8.4-thm:indef}: if $g$ is a differentiable one-to-one function from an open interval $J$ onto an open interval $I$ with $g'\ne0$, and $f$ is continuous on $I$, then $\int f(x)\,dx=\left[\int f(g(\theta))g'(\theta)\,d\theta\right]_{\theta=g^{-1}(x)}$.

(1) Put $g(u)=\frac{u^n-b}a$ on $(0,\infty)$. Then $g'(u)=\frac na\,u^{n-1}\ne0$ for $u>0$, so $g$ is strictly monotone (hence one-to-one) and continuous; its image is the open interval $\{x:ax+b>0\}$, and $g^{-1}(x)=(ax+b)^{1/n}$, because $u^n=ax+b$ and $u>0$. Take $J=g^{-1}(I)$. For $f(x)=R\left(x,(ax+b)^{1/n}\right)$ we get $f(g(u))=R\left(g(u),(u^n)^{1/n}\right)=R\left(\frac{u^n-b}a,u\right)$ since $u>0$. The theorem gives the formula; the integrand $R\left(\frac{u^n-b}a,u\right)\frac na u^{n-1}$ is a rational function of $u$ because $R$ is rational and the other factors are polynomials.

(2) Put $g(u)=\ln u$ on $(0,\infty)$: it is differentiable, one-to-one, onto $(-\infty,\infty)$, with $g'(u)=\frac1u\ne0$ and $g^{-1}(x)=e^x$. Here $f(x)=R(e^x)$, so $f(g(u))=R(e^{\ln u})=R(u)$, and the theorem gives $\int R(e^x)\,dx=\left[\int R(u)\frac1u\,du\right]_{u=e^x}$.
\end{proof}

\begin{exmp}{A root removed by a substitution}{m7}
Evaluate $\displaystyle\int\frac{x}{2+\sqrt{x+1}}\,dx$ on $x>-1$.

\textbf{Solution.} The only root is $\sqrt{x+1}$, so let $u=\sqrt{x+1}$ (Theorem~\ref{thm:rational}(1) with $n=2$, $a=1$, $b=1$). Then $x=u^2-1$, $dx=2u\,du$, and
\[\int\frac{x}{2+\sqrt{x+1}}\,dx=\int\frac{u^2-1}{2+u}\cdot2u\,du=2\int\frac{u^3-u}{u+2}\,du .\]
The integrand is an improper rational function, so divide: $u^3-u=(u+2)(u^2-2u+3)-6$ (check: $(u+2)(u^2-2u+3)=u^3-2u^2+3u+2u^2-4u+6=u^3-u+6$). Hence $\frac{u^3-u}{u+2}=u^2-2u+3-\frac6{u+2}$ and
\begin{align*}2\int\left(u^2-2u+3-\frac6{u+2}\right)du&=2\left(\frac{u^3}3-u^2+3u-6\ln(u+2)\right)\\&=\frac23u^3-2u^2+6u-12\ln(u+2).\end{align*}
Return to $x$ with $u=\sqrt{x+1}$, $u^2=x+1$, $u^3=(x+1)^{3/2}$:
\[\int\frac{x}{2+\sqrt{x+1}}\,dx=\frac23(x+1)^{3/2}-2(x+1)+6\sqrt{x+1}-12\ln\left(2+\sqrt{x+1}\right)+C .\]
%@ m7 | div(Symbol('u')**3-Symbol('u'),Symbol('u')+2,Symbol('u')) == (Symbol('u')**2-2*Symbol('u')+3,-6) and ANT(Q(2,3)*(x+1)**Q(3,2)-2*(x+1)+6*sqrt(x+1)-12*log(2+sqrt(x+1)),x/(2+sqrt(x+1)),[Q(1,2),Q(3,1),Q(8,1),-Q(1,2)]) == 0 | -12\ln\left(2+\sqrt{x+1}\right)+C .\]
\emph{Check.} With $u=\sqrt{x+1}$ we have $\frac{du}{dx}=\frac1{2u}$, so by the Chain Rule the derivative of the answer is $\frac1{2u}\left(2u^2-4u+6-\frac{12}{u+2}\right)$. Over a common denominator, $2u^2-4u+6-\frac{12}{u+2}=\frac{(2u^2-4u+6)(u+2)-12}{u+2}=\frac{2u^3-2u}{u+2}$, so the derivative is $\frac1{2u}\cdot\frac{2u(u^2-1)}{u+2}=\frac{u^2-1}{u+2}=\frac{x}{2+\sqrt{x+1}}$, the integrand.
\end{exmp}

\begin{kc}{Which substitution?}{kc2}
Find a substitution that turns each integral into one of rational type (do not evaluate): (a) $\displaystyle\int\frac{dx}{1+\sqrt{3x+2}}$; (b) $\displaystyle\int\frac{e^{2x}}{e^x+1}\,dx$; (c) $\displaystyle\int\frac{\sqrt x}{1+\sqrt[3]x}\,dx$.
\end{kc}

% ---------------------------------------------------------------------
\subsection*{Patterns: roots, rational functions, definite integrals}

\begin{exmp}{Do not overlook a simpler solution}{m8}
Evaluate $\displaystyle\int\frac{x^3}{\sqrt{4-x^2}}\,dx$ on $\lvert x\rvert<2$ in two ways.

\textbf{Solution.} \emph{Route A: substitution.} The root contains $4-x^2$, whose derivative $-2x$ shares a factor $x$ with the numerator. Write $x^3\,dx=x^2\cdot x\,dx$ and let $u=4-x^2$, so $x^2=4-u$ and $x\,dx=-\frac12\,du$:
\begin{align*}\int\frac{x^3\,dx}{\sqrt{4-x^2}}&=-\frac12\int\frac{4-u}{\sqrt u}\,du=-\frac12\int\left(4u^{-1/2}-u^{1/2}\right)du\\&=-\frac12\left(8u^{1/2}-\frac23u^{3/2}\right)=-4\sqrt u+\frac13u^{3/2}.\end{align*}
So the integral is $-4\sqrt{4-x^2}+\frac13(4-x^2)^{3/2}+C$.
%@ m8a | ANT(-4*sqrt(4-x**2)+(4-x**2)**Q(3,2)/3,x**3/sqrt(4-x**2),[Q(1,2),Q(1,1),Q(3,2),-Q(1,1)]) == 0 | So the integral is $-4\sqrt{4-x^2}+\frac13(4-x^2)^{3/2}+C$.

\emph{Route B: trigonometric substitution.} The pattern $\sqrt{a^2-x^2}$ with $a=2$ suggests $x=2\sin\theta$, $dx=2\cos\theta\,d\theta$, $\sqrt{4-x^2}=2\cos\theta$ (as $\cos\theta>0$ for $-\frac\pi2<\theta<\frac\pi2$). Then $\frac{x^3\,dx}{\sqrt{4-x^2}}=\frac{8\sin^3\theta\cdot2\cos\theta\,d\theta}{2\cos\theta}=8\sin^3\theta\,d\theta$. This is an odd power of sine (Section~8.3): $8\int(1-\cos^2\theta)\sin\theta\,d\theta=8\left(-\cos\theta+\frac{\cos^3\theta}3\right)$. With $\cos\theta=\frac{\sqrt{4-x^2}}2$ (Theorem~\ref{8.4-thm:back}) this is $-4\sqrt{4-x^2}+\frac83\cdot\frac{(4-x^2)^{3/2}}8$, which is the same function as in Route A.
%@ m8b | sum(NUM(-4*sqrt(4-v**2)+Q(8,3)*(sqrt(4-v**2)/2)**3,-4*sqrt(4-v**2)+(4-v**2)**Q(3,2)/3) for v in (Q(1,2),Q(1,1),Q(3,2))) == 0 | which is the same function as in Route A.

Both routes are correct; Route A is shorter because it uses no inverse function and no second technique. The pattern $\sqrt{a^2-x^2}$ ``suggests'' a trigonometric substitution, but an odd power of $x$ outside the root is the signal to look for $u=a^2-x^2$ first.
\end{exmp}

\begin{exmp}{A rational function: test, then decompose}{m9}
Evaluate $\displaystyle\int\frac{5x^2-x+7}{(x-1)(x^2+4)}\,dx$.

\textbf{Solution.} \emph{Tests first.} The fraction is proper (degree $2<3$). The derivative of the denominator $x^3-x^2+4x-4$ is $3x^2-2x+4$, which is not a constant multiple of the numerator, so a substitution $u=$ denominator fails. The factor $x^2+4$ is irreducible ($0-16<0$), so partial fractions:
\[\frac{5x^2-x+7}{(x-1)(x^2+4)}=\frac A{x-1}+\frac{Bx+C}{x^2+4},\qquad 5x^2-x+7=A(x^2+4)+(Bx+C)(x-1).\]
Cover up $x-1$ at $x=1$: $A=\frac{5-1+7}{1+4}=\frac{11}5$. Compare the $x^2$ coefficients: $5=A+B$, so $B=\frac{14}5$. Compare the constants: $7=4A-C=\frac{44}5-C$, so $C=\frac95$. Check the coefficient of $x$: $-1=C-B=\frac95-\frac{14}5=-1$. Now split the quadratic piece into a logarithm part and an inverse tangent part:
\[\frac{14x+9}{5(x^2+4)}=\frac75\cdot\frac{2x}{x^2+4}+\frac95\cdot\frac1{x^2+4},\]
and $\int\frac{dx}{x^2+4}=\frac12\tan^{-1}\frac x2$ (formula (11) with $a=2$). Therefore
\[\int\frac{5x^2-x+7}{(x-1)(x^2+4)}\,dx=\frac{11}5\ln\lvert x-1\rvert+\frac75\ln\left(x^2+4\right)+\frac9{10}\tan^{-1}\frac x2+C .\]
%@ m9 | simplify(apart((5*x**2-x+7)/((x-1)*(x**2+4)),x)-(Q(11,5)/(x-1)+(14*x+9)/(5*(x**2+4)))) == 0 and Q(5-1+7,1+4) == Q(11,5) and ANT(Q(11,5)*log(sqrt((x-1)**2))+Q(7,5)*log(x**2+4)+Q(9,10)*atan(x/2),(5*x**2-x+7)/((x-1)*(x**2+4)),[-Q(2,1),Q(1,2),Q(3,1),Q(5,1)]) == 0 | +\frac9{10}\tan^{-1}\frac x2+C .\]
\emph{Check.} Differentiate term by term. $\frac d{dx}\frac{11}5\ln\lvert x-1\rvert=\frac{11/5}{x-1}$; $\frac d{dx}\frac75\ln(x^2+4)=\frac{14x/5}{x^2+4}$; and by the Chain Rule $\frac d{dx}\frac9{10}\tan^{-1}\frac x2=\frac9{10}\cdot\frac{1/2}{1+x^2/4}=\frac{9/5}{x^2+4}$. The sum is $\frac{11/5}{x-1}+\frac{14x+9}{5(x^2+4)}$, the decomposition found above, which adds up to the integrand.
\end{exmp}

\begin{exmp}{A definite integral: substitution, then parts}{m10}
Evaluate $\displaystyle\int_0^{\pi^2/4}\sin\sqrt x\,dx$.

\textbf{Solution.} The only appearance of $x$ is inside the root, but $\frac{du}{dx}=\frac1{2\sqrt x}$ is not in the integrand. Instead solve for $x$: let $u=\sqrt x$, so $x=u^2$ and $dx=2u\,du$ (the idea of Theorem~\ref{thm:rational}, with an integrand that is not rational). The limits change: $x=0$ gives $u=0$ and $x=\frac{\pi^2}4$ gives $u=\frac\pi2$. Then
\[\int_0^{\pi^2/4}\sin\sqrt x\,dx=\int_0^{\pi/2}2u\sin u\,du .\]
This is a polynomial times a sine, so use parts: $p=2u$, $dv=\sin u\,du$, so $dp=2\,du$ and $v=-\cos u$ (the letter $p$ avoids a clash with the substitution variable):
\[\int_0^{\pi/2}2u\sin u\,du=\Big[-2u\cos u\Big]_0^{\pi/2}+2\int_0^{\pi/2}\cos u\,du=\left(-\pi\cdot0-0\right)+2\Big[\sin u\Big]_0^{\pi/2}=2 .\]
%@ m10 | NUM(Integral(sin(sqrt(x)),(x,0,pi**2/4)),2) == 0 and NUM(Integral(2*x*sin(x),(x,0,pi/2)),2) == 0 and sin(pi/2) == 1 | =\left(-\pi\cdot0-0\right)+2\Big[\sin u\Big]_0^{\pi/2}=2 .\]
\emph{Check.} Going back to $x$ (not to $u$): an antiderivative is $-2\sqrt x\cos\sqrt x+2\sin\sqrt x$. With $\frac d{dx}\sqrt x=\frac1{2\sqrt x}$ and the Chain Rule, $\frac d{dx}\left(-2\sqrt x\cos\sqrt x\right)=-\frac{\cos\sqrt x}{\sqrt x}+\sin\sqrt x$ and $\frac d{dx}\left(2\sin\sqrt x\right)=\frac{\cos\sqrt x}{\sqrt x}$, so the sum is $\sin\sqrt x$. At the limits: $-2\cdot\frac\pi2\cdot0+2\cdot1-(0+0)=2$.
%@ m10b | ANT(-2*sqrt(x)*cos(sqrt(x))+2*sin(sqrt(x)),sin(sqrt(x)),[Q(1,2),Q(2,1),Q(5,1)]) == 0 | At the limits: $-2\cdot\frac\pi2\cdot0+2\cdot1-(0+0)=2$.

\begin{center}
\begin{tikzpicture}
  \begin{scope}[x=2.1cm,y=1.9cm]
    \draw[->] (-0.1,0) -- (2.9,0) node[right] {$x$};
    \draw[->] (0,-0.05) -- (0,1.3) node[above] {$y$};
    \fill[blue!15] (0,0) -- plot[domain=0:2.4674,samples=60,variable=\t] ({\t},{sin(deg(sqrt(\t)))}) -- (2.4674,0) -- cycle;
    \draw[very thick,blue!70!black,domain=0:2.4674,samples=80,variable=\t] plot ({\t},{sin(deg(sqrt(\t)))});
    \fill (2.4674,1) circle (1.5pt) node[above right,font=\small] {$\left(\frac{\pi^2}4,1\right)$};
%@ f2-a | NUM(sin(sqrt(pi**2/4)),1) == 0 and abs(float(pi**2/4)-2.4674) < 1e-4 | \fill (2.4674,1) circle (1.5pt) node[above right,font=\small] {$\left(\frac{\pi^2}4,1\right)$};
    \node[below,font=\small] at (2.4674,0) {$\frac{\pi^2}4$};
    \node[font=\small] at (1.2,0.3) {area $2$};
    \node[font=\small\bfseries] at (1.3,-0.55) {(a) $\sin\sqrt x$ on $\left[0,\frac{\pi^2}4\right]$};
  \end{scope}
  \begin{scope}[xshift=7.6cm,x=2.0cm,y=0.95cm]
    \draw[->] (-0.1,0) -- (2.0,0) node[right] {$u$};
    \draw[->] (0,-0.1) -- (0,3.6) node[above] {$y$};
    \fill[blue!15] (0,0) -- plot[domain=0:1.5708,samples=60,variable=\t] ({\t},{2*\t*sin(deg(\t))}) -- (1.5708,0) -- cycle;
    \draw[very thick,blue!70!black,domain=0:1.5708,samples=60,variable=\t] plot ({\t},{2*\t*sin(deg(\t))});
    \fill (1.5708,3.1416) circle (1.5pt) node[above left,font=\small] {$\left(\frac\pi2,\pi\right)$};
%@ f2-b | NUM(2*(pi/2)*sin(pi/2),pi) == 0 and abs(float(pi/2)-1.5708) < 1e-4 | \fill (1.5708,3.1416) circle (1.5pt) node[above left,font=\small] {$\left(\frac\pi2,\pi\right)$};
    \node[below,font=\small] at (1.5708,0) {$\frac\pi2$};
    \node[font=\small] at (1.0,1.0) {area $2$};
    \node[font=\small\bfseries] at (1.0,-1.1) {(b) $2u\sin u$ on $\left[0,\frac\pi2\right]$};
  \end{scope}
\end{tikzpicture}
\end{center}
\begin{center}\small Figure 2. (a) The region under $y=\sin\sqrt x$ for $0\le x\le\frac{\pi^2}4$. (b) The region under $y=2u\sin u$ for $0\le u\le\frac\pi2$ after the substitution $x=u^2$. The regions have different shapes and different vertical scales ($1.9$ cm per unit in (a), $0.95$ cm per unit in (b)), but the same area, $2$.\end{center}
\end{exmp}

\begin{exmp}{A challenging rewrite}{m11}
Evaluate $\displaystyle\int\frac{\sin x}{\sin x+\cos x}\,dx$ on an interval where $\sin x+\cos x\ne0$.

\textbf{Solution.} Nothing in the table fits, no identity simplifies the quotient, and a substitution $u=\sin x+\cos x$ needs $du=(\cos x-\sin x)\,dx$, which is not the numerator. But the denominator and its derivative are related to the numerator: with $s=\sin x$ and $c=\cos x$ we have $2s=(s+c)+(s-c)$, so
\[\frac{\sin x}{\sin x+\cos x}=\frac12\cdot\frac{(s+c)+(s-c)}{s+c}=\frac12+\frac12\cdot\frac{\sin x-\cos x}{\sin x+\cos x}.\]
In the last fraction the numerator is \emph{minus} the derivative of the denominator. Let $u=\sin x+\cos x$, $du=(\cos x-\sin x)\,dx$, so $(\sin x-\cos x)\,dx=-du$:
\[\int\frac{\sin x}{\sin x+\cos x}\,dx=\frac12\int dx-\frac12\int\frac{du}u=\frac x2-\frac12\ln\lvert\sin x+\cos x\rvert+C .\]
%@ m11 | ANT(x/2-log(sqrt((sin(x)+cos(x))**2))/2,sin(x)/(sin(x)+cos(x)),[Q(1,2),Q(1,1),Q(2,1),-Q(1,3)]) == 0 | =\frac x2-\frac12\ln\lvert\sin x+\cos x\rvert+C .\]
\emph{Check.} $\frac d{dx}\left[\frac x2-\frac12\ln\lvert\sin x+\cos x\rvert\right]=\frac12-\frac12\cdot\frac{\cos x-\sin x}{\sin x+\cos x}=\frac{(\sin x+\cos x)-(\cos x-\sin x)}{2(\sin x+\cos x)}=\frac{2\sin x}{2(\sin x+\cos x)}=\frac{\sin x}{\sin x+\cos x}$.
\emph{A definite integral.} On $\left[0,\frac\pi2\right]$ we have $\sin x+\cos x\ge1>0$, and $\sin 0+\cos0=1=\sin\frac\pi2+\cos\frac\pi2$, so the logarithm terms cancel and $\int_0^{\pi/2}\frac{\sin x}{\sin x+\cos x}\,dx=\frac\pi4$. (Example~\ref{ex:b13} explains this value by symmetry.)
%@ m11b | NUM(Integral(sin(x)/(sin(x)+cos(x)),(x,0,pi/2)),pi/4) == 0 and log(sin(0)+cos(0)) == 0 | so the logarithm terms cancel and $\int_0^{\pi/2}\frac{\sin x}{\sin x+\cos x}\,dx=\frac\pi4$. (Example~\ref{ex:b13} explains this value by symmetry.)
\end{exmp}

\begin{think}{Two answers}{twoanswers}
Evaluate $\int\sin x\cos x\,dx$ with $u=\sin x$ and again with $u=\cos x$. The two answers look different. Are they both right? What exactly is the relation between them?
\end{think}

% ---------------------------------------------------------------------
\subsection*{Integrals with no elementary antiderivative}

Some integrals resist every method of this chapter, and that is not because we are not clever enough. First we must be clear about what is claimed.

\begin{defn}{Elementary function}{elem}
A function is \emph{elementary} if it can be built from constants and the function $x$ by finitely many sums, differences, products, quotients and compositions, and by taking roots, exponentials, logarithms, sines, cosines and their inverse functions. These are the functions that a scientific calculator can compute from its keys.
\end{defn}

Every antiderivative found in Sections 8.1 to 8.5 is elementary. But the statement \emph{``$\int e^{-x^2}\,dx$ is not elementary''} does \textbf{not} say that the integral does not exist. The function $e^{-x^2}$ is continuous, so by the Fundamental Theorem of Calculus the function $F(x)=\int_0^xe^{-t^2}\,dt$ is differentiable with $F'(x)=e^{-x^2}$: it \emph{is} an antiderivative. The point is that $F$ cannot be written with the functions of the definition. It is an honest new function, defined by an area; it is called the error function (up to a constant factor, $\operatorname{erf}x=\frac2{\sqrt\pi}F(x)$) and it appears in probability and statistics.

How can we know that no combination of elementary functions works? Trying many methods and failing proves nothing. What does prove it is a theorem, which we quote, together with an algebraic fact that we prove in full.

\begin{thm}{Liouville's criterion for $f\,e^g$ (quoted)}{liouville}
Let $f$ and $g$ be rational functions with $g$ not constant. If $\int f(x)\,e^{g(x)}\,dx$ has an elementary antiderivative on an interval, then there is a \emph{rational} function $R$ such that $R(x)e^{g(x)}$ is an antiderivative, that is,
\[R'+g'R=f .\]
\end{thm}
%%@omitted liouville
\emph{We do not prove Theorem~\ref{thm:liouville}.} It is a theorem of Liouville (1830s) whose proof uses differential algebra and is beyond this course; it is the one result of this section that is quoted and not proved, and no appendix section supplies it. The converse is easy and we have already used it: if $R'+g'R=f$ then $\frac d{dx}\left(Re^g\right)=\left(R'+g'R\right)e^g=fe^g$. The next theorem shows that for the integrands of interest the equation $R'+g'R=f$ has \emph{no} rational solution. Its proof uses only Euclid's lemma (Theorem~\ref{8.5-thm:euclid}) and the degrees of polynomials.

\begin{thm}{No rational solution}{norat}
\begin{enumerate}[label=(\alph*)]
  \item If $h$ is a nonconstant polynomial, there is no rational function $R$ with $R'+hR=1$.
  \item There is no rational function $R$ with $R'+R=\frac1x$.
\end{enumerate}
\end{thm}
\begin{proof}
Write $R=\frac PQ$ in lowest terms ($P$ and $Q$ polynomials with no common nonconstant factor). In case (a) put $s=1$ and in case (b) put $s=x$ and $h=1$. Multiplying the equation by $sQ^2$ turns it, in both cases, into
\[s\left(P'Q-PQ'+hPQ\right)=Q^2 .\tag{$*$}\]

\emph{Step 1: $Q$ is constant.} Suppose not. Let $p$ be an irreducible factor of $Q$ and write $Q=p^mT$ with $m\ge1$ and $p$ not dividing $T$. Then $Q'=mp^{m-1}p'T+p^mT'$, and
\[P'Q-PQ'+hPQ=p^{m-1}\left[p\,P'T-mPp'T-p\,PT'+p\,hPT\right]=p^{m-1}\,W,\qquad W=p\,U-mPp'T,\]
where $U=P'T-PT'+hPT$. We claim that $p$ does not divide $W$. If it did, $p$ would divide $mPp'T$. But $p$ does not divide $P$ (otherwise $P$ and $Q$ would have the common factor $p$), $p$ does not divide $T$ (choice of $m$), $p$ does not divide $m$ (a nonzero constant), and $p$ does not divide $p'$ (the polynomial $p'$ is nonzero with degree less than $\deg p$). By Euclid's lemma (Theorem~\ref{8.5-thm:euclid}) applied repeatedly, $p$ does not divide the product $mPp'T$, a contradiction. So $p\nmid W$.

Now $(*)$ reads $s\,p^{m-1}W=p^{2m}T^2$, and cancelling $p^{m-1}$ gives $sW=p^{m+1}T^2$. Hence $p$ divides $sW$. If $s=1$ then $p$ divides $W$, which is false. If $s=x$, Euclid's lemma gives $p\mid x$ or $p\mid W$; the second is false, so $p$ divides $x$, and since $p$ is not constant, $p=cx$ for a constant $c\ne0$. Then $x=p/c$, and $sW=p^{m+1}T^2$ gives $W=c\,p^mT^2$, which is divisible by $p$ (because $m\ge1$): again a contradiction. So $Q$ is constant, and $R$ is a polynomial (absorb the constant into $P$): $Q=1$.

\emph{Step 2: $R=P$ is not a polynomial solution.} (a) If $P=0$ the equation says $0=1$. If $P\ne0$, then $\deg(hP)=\deg h+\deg P\ge\deg P+1>\deg P'$, so $\deg\left(P'+hP\right)=\deg h+\deg P\ge1$, but the right side $1$ has degree $0$. (b) The equation $P'+P=\frac1x$ gives $x\left(P'+P\right)=1$; at $x=0$ the left side is $0$ and the right side is $1$. In both cases there is no solution.
\end{proof}

\begin{thm}{Consequences}{conseq}
The functions $e^{x^2}$, $e^{-x^2}$, $e^{x^3}$, $\frac{e^x}x$ and $\frac1{\ln x}$ (for $x>1$) have no elementary antiderivative.
\end{thm}
\begin{proof}
For $e^{g}$ with $g=x^2,-x^2,x^3$ take $f=1$ in Theorem~\ref{thm:liouville}: an elementary antiderivative would give a rational $R$ with $R'+g'R=1$, where $g'=2x,-2x,3x^2$ is a nonconstant polynomial; this contradicts Theorem~\ref{thm:norat}(a). For $\frac{e^x}x$ take $g=x$ and $f=\frac1x$: Theorem~\ref{thm:norat}(b). For $\frac1{\ln x}$: if $G$ were an elementary antiderivative on $x>1$, then $G(e^u)$ (an elementary function of $u$) would satisfy $\frac d{du}G(e^u)=G'(e^u)\,e^u=\frac{e^u}{\ln e^u}=\frac{e^u}u$ for $u>0$, an elementary antiderivative of $\frac{e^u}u$, which does not exist.
\end{proof}

\begin{proc}{Recognizing a non-elementary integral}{nonelem}
\begin{enumerate}[nosep]
  \item \emph{Look for the warning sign:} a transcendental function ($e^{(\cdot)}$, $\sin$, $\cos$, $\ln$) of a \emph{non-linear} argument ($x^2$, $x^3$, $\ln x$, $\frac1x$) whose derivative is \emph{not} present. Compare $\int xe^{-x^2}\,dx$ (the derivative $-2x$ is there, so substitution works) with $\int e^{-x^2}\,dx$ (it is not).
  \item \emph{Watch the methods.} If parts returns the same kind of integral (or a harder one) and a substitution leaves an $x$ behind, suspect that no elementary antiderivative exists.
  \item \emph{Test with the criterion.} For $\int f\,e^g\,dx$ with $f,g$ rational, look for a rational $R$ with $R'+g'R=f$ (Theorems~\ref{thm:liouville} and~\ref{thm:norat}).
  \item \emph{Know the standard list} (each follows from Liouville's theorem and its extensions, proofs omitted except those above): $\int e^{\pm x^2}dx$, $\int\frac{e^x}x\,dx$, $\int\frac{dx}{\ln x}$, $\int\frac{\sin x}x\,dx$, $\int\sin\left(x^2\right)dx$, $\int\cos\left(x^2\right)dx$, $\int\sqrt{1+x^3}\,dx$.
  \item \emph{A failed search is not a proof}, and a computer's inability to find a formula is not one either. But a proved non-elementary integral is not the end: \emph{definite} integrals still have values, found numerically or bracketed by inequalities, and $\int_0^xf(t)\,dt$ is a perfectly good function.
\end{enumerate}
\end{proc}

\begin{exmp}{Elementary or not?}{m12}
Decide whether each integral has an elementary antiderivative; find it, or give the numerical value and a rough bracket.

(a) $\displaystyle\int xe^{-x^2}\,dx$. \quad (b) $\displaystyle\int e^{-x^2}\,dx$. \quad (c) $\displaystyle\int_0^1e^{-x^2}\,dx$.

\textbf{Solution.} (a) The exponent $-x^2$ has derivative $-2x$, and $x$ is present up to the factor $-2$. Let $u=-x^2$, $du=-2x\,dx$, $x\,dx=-\frac12\,du$:
\[\int xe^{-x^2}\,dx=-\frac12\int e^u\,du=-\frac12e^{-x^2}+C .\]
%@ m12a | ANT(-exp(-x**2)/2,x*exp(-x**2),[Q(1,2),Q(1,1),-Q(3,2),Q(2,1)]) == 0 | =-\frac12e^{-x^2}+C .\]
\emph{Check.} $\frac d{dx}\left(-\frac12e^{-x^2}\right)=-\frac12e^{-x^2}\cdot(-2x)=xe^{-x^2}$.

(b) Here the derivative $-2x$ is missing. Parts with $u=e^{-x^2}$, $dv=dx$ gives $xe^{-x^2}+2\int x^2e^{-x^2}\,dx$, and parts on the new integral (with $p=x$, $dv=xe^{-x^2}\,dx$) brings back exactly $\int e^{-x^2}\,dx$: the two steps together give the empty identity $I=I$. This is the warning sign. By Theorem~\ref{thm:conseq}, $e^{-x^2}$ has \emph{no} elementary antiderivative. It does have the antiderivative $F(x)=\int_0^xe^{-t^2}\,dt=\frac{\sqrt\pi}2\operatorname{erf}x$ (Figure~3), with $F'(x)=e^{-x^2}$, $F(1)\approx0.7468$, and, from a later section on improper integrals, $F(x)\to\frac{\sqrt\pi}2\approx0.8862$ as $x\to\infty$.
%@ m12b | ANT(sqrt(pi)*erf(x)/2,exp(-x**2),[Q(1,2),Q(1,1),-Q(3,2),Q(2,1)]) == 0 and integrate(exp(-x**2),(x,0,oo)) == sqrt(pi)/2 | $F(x)\to\frac{\sqrt\pi}2\approx0.8862$ as $x\to\infty$.

(c) By the Fundamental Theorem, $\int_0^1e^{-x^2}\,dx=F(1)\approx0.7468$; no antiderivative formula is needed. We can also bracket it by hand with two inequalities from Calculus~I, valid for all $t$: $e^{-t}\ge1-t$ (the graph of $e^{-t}$ lies above its tangent line at $t=0$) and $e^t\ge1+t$ for $t\ge0$, hence $e^{-t}\le\frac1{1+t}$. With $t=x^2$, for $0\le x\le1$,
\[1-x^2\le e^{-x^2}\le\frac1{1+x^2},\qquad\text{so}\qquad\int_0^1\left(1-x^2\right)dx=\frac23\le\int_0^1e^{-x^2}\,dx\le\int_0^1\frac{dx}{1+x^2}=\frac\pi4 .\]
Numerically $0.6667\le0.7468\le0.7854$.
%@ m12c | NUM(Integral(exp(-x**2),(x,0,1)),sqrt(pi)*erf(1)/2) == 0 and Q(2,3) < Integral(exp(-x**2),(x,0,1)).evalf(30) and Integral(exp(-x**2),(x,0,1)).evalf(30) < (pi/4).evalf(30) and abs(float(sqrt(pi)*erf(1)/2)-0.7468) < 1e-4 and integrate(1-x**2,(x,0,1)) == Q(2,3) and integrate(1/(1+x**2),(x,0,1)) == pi/4 | Numerically $0.6667\le0.7468\le0.7854$.

\begin{center}
\begin{tikzpicture}
  \begin{scope}[x=1.6cm,y=2.8cm]
    \draw[->] (-0.4,0) -- (2.7,0) node[right] {$x$};
    \draw[->] (0,-0.05) -- (0,1.25) node[above] {$y$};
    \fill[blue!15] (0,0) -- plot[domain=0:1,samples=40,variable=\t] ({\t},{exp(-\t*\t)}) -- (1,0) -- cycle;
    \draw[very thick,blue!70!black,domain=-0.4:2.6,samples=80,variable=\t] plot ({\t},{exp(-\t*\t)});
    \fill (0,1) circle (1.5pt) node[above left,font=\small] {$(0,1)$};
%@ f3-a | exp(-0**2) == 1 | \fill (0,1) circle (1.5pt) node[above left,font=\small] {$(0,1)$};
    \fill (1,0.3679) circle (1.5pt) node[above right,font=\small] {$(1,e^{-1})$};
%@ f3-b | abs(float(exp(-1))-0.3679) < 1e-4 and exp(-1**2) == exp(-1) | \fill (1,0.3679) circle (1.5pt) node[above right,font=\small] {$(1,e^{-1})$};
    \node[below,font=\small] at (1,0) {$1$};
    \node[font=\small] at (0.5,0.22) {area $\approx0.7468$};
    \node[font=\small\bfseries] at (1.1,-0.4) {(a) $y=e^{-x^2}$};
  \end{scope}
  \begin{scope}[xshift=7.4cm,x=1.4cm,y=3.2cm]
    \draw[->] (-0.2,0) -- (3.3,0) node[right] {$x$};
    \draw[->] (0,-0.05) -- (0,1.05) node[above] {$y$};
    \draw[dashed,gray] (0,0.8862) -- (3.2,0.8862);
    \node[left,font=\small] at (0,0.8862) {$\frac{\sqrt\pi}2$};
    \draw[very thick,green!50!black] plot[smooth] coordinates {(0,0) (0.25,0.2449) (0.5,0.4613) (0.75,0.6302) (1,0.7468) (1.25,0.8179) (1.5,0.8562) (1.75,0.8744) (2,0.8821) (2.25,0.8849) (2.5,0.8859) (2.75,0.8861) (3,0.8862)};
%@ f3-c | (lambda L: all(abs(float(sqrt(pi)*erf(Q(i,4))/2)-L[i]) < 1e-4 for i in range(len(L))))((0,0.2449,0.4613,0.6302,0.7468,0.8179,0.8562,0.8744,0.8821,0.8849,0.8859,0.8861,0.8862)) | \draw[very thick,green!50!black] plot[smooth] coordinates {(0,0) (0.25,0.2449) (0.5,0.4613) (0.75,0.6302) (1,0.7468) (1.25,0.8179) (1.5,0.8562) (1.75,0.8744) (2,0.8821) (2.25,0.8849) (2.5,0.8859) (2.75,0.8861) (3,0.8862)};
    \fill (1,0.7468) circle (1.5pt) node[below right,font=\small] {$(1,F(1))$};
%@ f3-d | NUM(Integral(exp(-x**2),(x,0,1)),sqrt(pi)*erf(1)/2) == 0 and abs(float(sqrt(pi)/2)-0.8862) < 1e-4 | \fill (1,0.7468) circle (1.5pt) node[below right,font=\small] {$(1,F(1))$};
    \node[below,font=\small] at (1,0) {$1$};
    \node[font=\small\bfseries] at (1.5,-0.3) {(b) $F(x)=\int_0^xe^{-t^2}\,dt$};
  \end{scope}
\end{tikzpicture}
\end{center}
\begin{center}\small Figure 3. (a) The graph of $e^{-x^2}$; the shaded area from $0$ to $1$ is $F(1)\approx0.7468$. (b) The area function $F(x)=\int_0^xe^{-t^2}\,dt$: it is an antiderivative of $e^{-x^2}$, increasing and levelling off at $\frac{\sqrt\pi}2\approx0.886$ (dashed), although no elementary formula for it exists. The vertical scales are $2.8$ cm and $3.2$ cm per unit.\end{center}
\end{exmp}

\begin{kc}{Elementary or not?}{kc3}
Use the procedure above and the theorems to decide whether each integral has an elementary antiderivative, and say why: (a) $\displaystyle\int xe^{x^2}\,dx$; (b) $\displaystyle\int e^{x^2}\,dx$; (c) $\displaystyle\int\cos x\,e^{\sin x}\,dx$; (d) $\displaystyle\int\frac{e^x}x\,dx$.
\end{kc}

\begin{think}{Failing is not proving}{provenot}
A student tries parts, three substitutions and a table on $\int e^{-x^2}\,dx$, fails every time, and writes ``so it has no elementary antiderivative''. Why is this not a proof? Which parts of this section do prove the claim, and which part of that proof did we not prove here?
\end{think}

\subsection*{Looking ahead}

When an integral cannot be done in closed form, or when the closed form is too long to be useful, the value of a \emph{definite} integral can still be approximated to any accuracy. The next sections introduce the use of tables and computer algebra (a table entry is a claim to be checked by differentiation, as above), numerical integration with error bounds, and improper integrals, which will give the value $\int_0^\infty e^{-x^2}\,dx=\frac{\sqrt\pi}2$ quoted in Example~\ref{ex:m12}. In the chapters on series, a function such as $\int_0^xe^{-t^2}\,dt$ will be computed by integrating the series of $e^{-t^2}$ term by term; that is how calculators evaluate it.

% =====================================================================
\subsection*{Example bank}

\textbf{What's in this bank.} Fourteen further worked examples, grouped by the pattern in the integrand and ordered from easiest to hardest within each group.
\begin{itemize}[nosep]
  \item Group 1: rewrite the integrand first (expand and split, multiply by a conjugate).
  \item Group 2: substitution, including a substitution that is not obvious and two substitutions in a row.
  \item Group 3: products (a cyclic integral, and a substitution followed by parts).
  \item Group 4: trigonometric integrands (odd powers, tangent and secant).
  \item Group 5: roots and quadratics (a trigonometric substitution, and a case where a plain substitution is better).
  \item Group 6: rational integrands and roots (a substitution that beats partial fractions, a sixth-root substitution).
  \item Group 7: proofs (a symmetry trick for definite integrals, a reduction formula).
\end{itemize}

\subsection*{Group 1: Rewrite first}

\begin{exmp}{Expand and split}{b1}
Evaluate $\displaystyle\int\frac{\left(\sqrt x+3\right)^2}x\,dx$ on $x>0$.

\textbf{Solution.} Expand the square: $\left(\sqrt x+3\right)^2=x+6\sqrt x+9$. Divide each term by $x$ (a numerator may be split): $\frac{x}x+\frac{6\sqrt x}x+\frac9x=1+6x^{-1/2}+\frac9x$. Integrate each term with the power rule and $\int\frac{dx}x=\ln x$ (for $x>0$):
\[\int\left(1+6x^{-1/2}+\frac9x\right)dx=x+6\cdot\frac{x^{1/2}}{1/2}+9\ln x+C=x+12\sqrt x+9\ln x+C .\]
%@ b1 | ANT(x+12*sqrt(x)+9*log(x),(sqrt(x)+3)**2/x,[Q(1,2),Q(1,1),Q(4,1),Q(9,1)]) == 0 | =x+12\sqrt x+9\ln x+C .\]
\emph{Check.} $\frac d{dx}\left(x+12\sqrt x+9\ln x\right)=1+\frac6{\sqrt x}+\frac9x=\frac{x+6\sqrt x+9}x=\frac{(\sqrt x+3)^2}x$.
\end{exmp}

\begin{exmp}{Multiply by a conjugate}{b2}
Evaluate $\displaystyle\int\frac{dx}{1+\sin x}$ on an interval where $\cos x\ne0$.

\textbf{Solution.} Nothing in the table fits, but $1-\sin^2x=\cos^2x$ makes a product of conjugates useful. Multiply the numerator and denominator by $1-\sin x$:
\[\frac1{1+\sin x}=\frac{1-\sin x}{(1+\sin x)(1-\sin x)}=\frac{1-\sin x}{1-\sin^2x}=\frac{1-\sin x}{\cos^2x}=\sec^2x-\sec x\tan x .\]
Both terms are in the table (formulas (5) and (7)):
\[\int\frac{dx}{1+\sin x}=\tan x-\sec x+C .\]
%@ b2 | ANT(tan(x)-sec(x),1/(1+sin(x)),[Q(1,2),Q(1,1),Q(3,2),-Q(1,1)]) == 0 | =\tan x-\sec x+C .\]
\emph{Check.} $\frac d{dx}(\tan x-\sec x)=\sec^2x-\sec x\tan x=\frac{1-\sin x}{\cos^2x}=\frac{1-\sin x}{(1-\sin x)(1+\sin x)}=\frac1{1+\sin x}$.
\end{exmp}

\subsection*{Group 2: Substitution}

\begin{exmp}{An exponential inside a quadratic}{b3}
Evaluate $\displaystyle\int\frac{e^x}{e^{2x}+4e^x+13}\,dx$.

\textbf{Solution.} All the $x$ appears in $e^x$, and $e^x\,dx$ is there. Let $u=e^x$, $du=e^x\,dx$, $e^{2x}=u^2$:
\[\int\frac{e^x\,dx}{e^{2x}+4e^x+13}=\int\frac{du}{u^2+4u+13} .\]
Complete the square: $u^2+4u+13=(u+2)^2+9$ (since $13-4=9>0$, the quadratic has no real root). By formula (11) with $a=3$ and the shift $u+2$ (Theorem~\ref{8.1-thm:quad}(i)),
\[\int\frac{du}{(u+2)^2+9}=\frac13\tan^{-1}\frac{u+2}3+C=\frac13\tan^{-1}\frac{e^x+2}3+C .\]
%@ b3 | expand((x+2)**2+9-(x**2+4*x+13)) == 0 and ANT(atan((exp(x)+2)/3)/3,exp(x)/(exp(2*x)+4*exp(x)+13),[0,Q(1,2),Q(1,1),-Q(1,1),Q(2,1)]) == 0 | =\frac13\tan^{-1}\frac{u+2}3+C=\frac13\tan^{-1}\frac{e^x+2}3+C .\]
\emph{Check.} $\frac d{dx}\frac13\tan^{-1}\frac{e^x+2}3=\frac13\cdot\frac{e^x/3}{1+(e^x+2)^2/9}=\frac{e^x}{9+(e^x+2)^2}=\frac{e^x}{e^{2x}+4e^x+13}$.
\end{exmp}

\begin{exmp}{Two substitutions in a row}{b4}
Evaluate $\displaystyle\int\frac{\ln(\ln x)}{x\ln x}\,dx$ on $x>1$.

\textbf{Solution.} The factor $\frac1x$ is the derivative of $\ln x$, so let $u=\ln x$, $du=\frac{dx}x$ (here $u>0$):
\[\int\frac{\ln(\ln x)}{x\ln x}\,dx=\int\frac{\ln u}{u}\,du .\]
Again the derivative of $\ln u$, which is $\frac1u$, is present. Let $w=\ln u$, $dw=\frac{du}u$:
\[\int\frac{\ln u}u\,du=\int w\,dw=\frac{w^2}2+C=\frac12\left(\ln u\right)^2+C=\frac12\left(\ln\ln x\right)^2+C .\]
%@ b4 | ANT(log(log(x))**2/2,log(log(x))/(x*log(x)),[Q(3,1),Q(5,1),Q(10,1),Q(3,2)]) == 0 | =\frac12\left(\ln u\right)^2+C=\frac12\left(\ln\ln x\right)^2+C .\]
\emph{Check.} $\frac d{dx}\frac12(\ln\ln x)^2=\ln\ln x\cdot\frac1{\ln x}\cdot\frac1x$.
\end{exmp}

\subsection*{Group 3: Products}

\begin{exmp}{A cyclic integral}{b5}
Evaluate $\displaystyle\int e^{2x}\cos3x\,dx$.

\textbf{Solution.} Let $I=\int e^{2x}\cos3x\,dx$ and $J=\int e^{2x}\sin3x\,dx$. \emph{First parts:} $p=\cos3x$, $dv=e^{2x}dx$, so $dp=-3\sin3x\,dx$ and $v=\frac12e^{2x}$:
\[I=\frac12e^{2x}\cos3x+\frac32J .\]
\emph{Second parts:} $p=\sin3x$, $dv=e^{2x}dx$, $dp=3\cos3x\,dx$, $v=\frac12e^{2x}$:
\[J=\frac12e^{2x}\sin3x-\frac32I .\]
Substitute $J$ into the first equation: $I=\frac12e^{2x}\cos3x+\frac34e^{2x}\sin3x-\frac94I$. Collect the $I$'s: $\frac{13}4I=e^{2x}\left(\frac12\cos3x+\frac34\sin3x\right)$, so
\[\int e^{2x}\cos3x\,dx=\frac{e^{2x}\left(2\cos3x+3\sin3x\right)}{13}+C ,\]
in agreement with Theorem~\ref{8.2-thm:cyc} ($a=2$, $b=3$).
%@ b5 | Q(1,1)+Q(9,4) == Q(13,4) and ANT(exp(2*x)*(2*cos(3*x)+3*sin(3*x))/13,exp(2*x)*cos(3*x),[Q(1,2),Q(1,1),Q(3,2),-Q(1,1)]) == 0 | in agreement with Theorem~\ref{8.2-thm:cyc} ($a=2$, $b=3$).
\end{exmp}

\begin{exmp}{Substitution, then parts}{b6}
Evaluate $\displaystyle\int x^3e^{-x^2}\,dx$.

\textbf{Solution.} Parts directly would need $\int x^3\,dx$ times $e^{-x^2}$ and gets worse. But $x^3=x^2\cdot x$ and $x\,dx$ is (up to a factor) the derivative of the exponent. Let $u=x^2$, $du=2x\,dx$, $x^3\,dx=u\cdot\frac{du}2$:
\[\int x^3e^{-x^2}\,dx=\frac12\int ue^{-u}\,du .\]
Now parts with $p=u$, $dv=e^{-u}du$, $dp=du$, $v=-e^{-u}$: $\int ue^{-u}\,du=-ue^{-u}+\int e^{-u}\,du=-ue^{-u}-e^{-u}=-(u+1)e^{-u}$. Therefore
\[\int x^3e^{-x^2}\,dx=-\frac12\left(x^2+1\right)e^{-x^2}+C .\]
%@ b6 | ANT(-(x**2+1)*exp(-x**2)/2,x**3*exp(-x**2),[Q(1,2),Q(1,1),-Q(3,2),Q(2,1)]) == 0 | =-\frac12\left(x^2+1\right)e^{-x^2}+C .\]
\emph{Check.} $\frac d{dx}\left[-\frac12(x^2+1)e^{-x^2}\right]=-\frac12\left[2xe^{-x^2}-2x(x^2+1)e^{-x^2}\right]=-xe^{-x^2}\left[1-x^2-1\right]=x^3e^{-x^2}$.
\end{exmp}

\subsection*{Group 4: Trigonometric integrands}

\begin{exmp}{An odd power of sine over a power of cosine}{b7}
Evaluate $\displaystyle\int\frac{\sin^3x}{\cos^2x}\,dx$ on an interval where $\cos x\ne0$.

\textbf{Solution.} The power of sine is odd: save one $\sin x$ and write $\sin^2x=1-\cos^2x$. Let $u=\cos x$, $du=-\sin x\,dx$:
\[\int\frac{(1-\cos^2x)\sin x}{\cos^2x}\,dx=-\int\frac{1-u^2}{u^2}\,du=-\int\left(u^{-2}-1\right)du=u^{-1}+u+C=\sec x+\cos x+C .\]
%@ b7 | ANT(sec(x)+cos(x),sin(x)**3/cos(x)**2,[Q(1,2),Q(1,1),Q(3,2),-Q(1,1)]) == 0 | =u^{-1}+u+C=\sec x+\cos x+C .\]
\emph{Check.} $\frac d{dx}(\sec x+\cos x)=\sec x\tan x-\sin x=\frac{\sin x}{\cos^2x}-\sin x=\frac{\sin x(1-\cos^2x)}{\cos^2x}=\frac{\sin^3x}{\cos^2x}$.
\end{exmp}

\begin{exmp}{Tangent and secant, odd power of secant}{b8}
Evaluate $\displaystyle\int\tan^3x\sec^3x\,dx$ on an interval where $\cos x\ne0$.

\textbf{Solution.} The power of tangent is odd, so save $\sec x\tan x$ (the derivative of $\sec x$) and convert the remaining $\tan^2x=\sec^2x-1$ (Section~8.3):
\[\tan^3x\sec^3x=\left(\sec^2x-1\right)\sec^2x\cdot\sec x\tan x .\]
Let $u=\sec x$, $du=\sec x\tan x\,dx$:
\[\int\left(u^2-1\right)u^2\,du=\frac{u^5}5-\frac{u^3}3+C=\frac{\sec^5x}5-\frac{\sec^3x}3+C .\]
%@ b8 | ANT(sec(x)**5/5-sec(x)**3/3,tan(x)**3*sec(x)**3,[Q(1,2),Q(1,1),Q(3,2),-Q(1,1)]) == 0 | =\frac{\sec^5x}5-\frac{\sec^3x}3+C .\]
\emph{Check.} $\frac d{dx}\left(\frac{\sec^5x}5-\frac{\sec^3x}3\right)=\sec^5x\tan x-\sec^3x\tan x=\sec^3x\tan x\left(\sec^2x-1\right)=\sec^3x\tan^3x$.
\end{exmp}

\subsection*{Group 5: Roots and quadratics}

\begin{exmp}{A trigonometric substitution}{b9}
Evaluate $\displaystyle\int\frac{dx}{x^2\sqrt{x^2+9}}$ on $x\ne0$.

\textbf{Solution.} The root $\sqrt{x^2+9}$ has the pattern $\sqrt{a^2+x^2}$ with $a=3$, and no plain substitution is visible. Let $x=3\tan\theta$, $-\frac\pi2<\theta<\frac\pi2$, $dx=3\sec^2\theta\,d\theta$, $\sqrt{x^2+9}=3\sec\theta$:
\[\int\frac{3\sec^2\theta\,d\theta}{9\tan^2\theta\cdot3\sec\theta}=\frac19\int\frac{\sec\theta}{\tan^2\theta}\,d\theta=\frac19\int\frac{\cos\theta}{\sin^2\theta}\,d\theta,\]
since $\frac{\sec\theta}{\tan^2\theta}=\frac1{\cos\theta}\cdot\frac{\cos^2\theta}{\sin^2\theta}=\frac{\cos\theta}{\sin^2\theta}$. With $w=\sin\theta$, $dw=\cos\theta\,d\theta$, this is $\frac19\int w^{-2}\,dw=-\frac1{9\sin\theta}$. By Theorem~\ref{8.4-thm:back}, $\sin\theta=\frac x{\sqrt{x^2+9}}$, so
\[\int\frac{dx}{x^2\sqrt{x^2+9}}=-\frac{\sqrt{x^2+9}}{9x}+C\qquad(x\ne0).\]
%@ b9 | ANT(-sqrt(x**2+9)/(9*x),1/(x**2*sqrt(x**2+9)),[Q(1,2),Q(1,1),Q(3,1),-Q(2,1)]) == 0 | =-\frac{\sqrt{x^2+9}}{9x}+C\qquad(x\ne0).\]
\emph{Check.} $\frac d{dx}\left(-\frac{\sqrt{x^2+9}}{9x}\right)=-\frac19\cdot\frac{\frac{x}{\sqrt{x^2+9}}\cdot x-\sqrt{x^2+9}}{x^2}=-\frac19\cdot\frac{x^2-(x^2+9)}{x^2\sqrt{x^2+9}}=\frac{1}{x^2\sqrt{x^2+9}}$.
\end{exmp}

\begin{exmp}{A root, but a plain substitution is better}{b10}
Evaluate $\displaystyle\int x^3\sqrt{x^2+4}\,dx$.

\textbf{Solution.} The root again has the form $\sqrt{x^2+a^2}$, but an \emph{odd} power of $x$ stands outside, so try $u=x^2+4$ first: $du=2x\,dx$, $x^2=u-4$, $x^3\,dx=x^2\cdot x\,dx=(u-4)\frac{du}2$:
\begin{align*}\int x^3\sqrt{x^2+4}\,dx&=\frac12\int(u-4)\,u^{1/2}\,du=\frac12\int\left(u^{3/2}-4u^{1/2}\right)du\\&=\frac12\left(\frac25u^{5/2}-\frac83u^{3/2}\right)=\frac{u^{5/2}}5-\frac{4u^{3/2}}3 .\end{align*}
So the integral is $\frac15\left(x^2+4\right)^{5/2}-\frac43\left(x^2+4\right)^{3/2}+C$.
%@ b10 | ANT((x**2+4)**Q(5,2)/5-4*(x**2+4)**Q(3,2)/3,x**3*sqrt(x**2+4),[Q(1,2),Q(1,1),Q(3,2),-Q(1,1)]) == 0 | So the integral is $\frac15\left(x^2+4\right)^{5/2}-\frac43\left(x^2+4\right)^{3/2}+C$.
\emph{Check.} $\frac d{dx}\left[\frac15u^{5/2}-\frac43u^{3/2}\right]=\left(\frac12u^{3/2}-2u^{1/2}\right)\cdot2x=x\,u^{1/2}(u-4)=x\cdot x^2\sqrt{x^2+4}$.
\end{exmp}

\subsection*{Group 6: Rational integrands and roots}

\begin{exmp}{Substitution beats partial fractions}{b11}
Evaluate $\displaystyle\int\frac{2x^3+3x}{\left(x^2+1\right)^2}\,dx$.

\textbf{Solution.} A repeated irreducible quadratic would need four unknowns. But the numerator is $x(2x^2+3)$, with a factor $x$ (an odd power) in front of a function of $x^2$. Let $u=x^2+1$, $du=2x\,dx$, $x\,dx=\frac{du}2$, $2x^2+3=2(u-1)+3=2u+1$:
\begin{align*}\int\frac{x(2x^2+3)}{(x^2+1)^2}\,dx&=\int\frac{2u+1}{u^2}\cdot\frac{du}2=\int\left(\frac1u+\frac1{2u^2}\right)du\\&=\ln u-\frac1{2u}+C=\ln\left(x^2+1\right)-\frac1{2(x^2+1)}+C .\end{align*}
%@ b11 | ANT(log(x**2+1)-1/(2*(x**2+1)),(2*x**3+3*x)/(x**2+1)**2,[Q(1,2),Q(1,1),-Q(3,2),Q(2,1)]) == 0 | =\ln u-\frac1{2u}+C=\ln\left(x^2+1\right)-\frac1{2(x^2+1)}+C .\end{align*}
\emph{Check.} $\frac d{dx}\left[\ln(x^2+1)-\frac1{2(x^2+1)}\right]=\frac{2x}{x^2+1}+\frac{x}{(x^2+1)^2}=\frac{2x(x^2+1)+x}{(x^2+1)^2}=\frac{2x^3+3x}{(x^2+1)^2}$.
\end{exmp}

\begin{exmp}{Two roots: a sixth-root substitution}{b12}
Evaluate $\displaystyle\int\frac{dx}{\sqrt x+\sqrt[3]x}$ on $x>0$.

\textbf{Solution.} The roots have indices $2$ and $3$, so use the least common multiple $6$: let $x=u^6$ with $u=x^{1/6}>0$. Then $dx=6u^5\,du$, $\sqrt x=u^3$ and $\sqrt[3]x=u^2$:
\[\int\frac{6u^5\,du}{u^3+u^2}=6\int\frac{u^3}{u+1}\,du\qquad\left(\text{cancel }u^2\right).\]
Divide: $u^3=(u+1)(u^2-u+1)-1$ (check: $(u+1)(u^2-u+1)=u^3+1$), so $\frac{u^3}{u+1}=u^2-u+1-\frac1{u+1}$ and
\[6\int\left(u^2-u+1-\frac1{u+1}\right)du=2u^3-3u^2+6u-6\ln(u+1)+C .\]
Back to $x$: $u^3=\sqrt x$, $u^2=\sqrt[3]x$, $u=x^{1/6}$:
\[\int\frac{dx}{\sqrt x+\sqrt[3]x}=2\sqrt x-3\sqrt[3]x+6\sqrt[6]x-6\ln\left(1+\sqrt[6]x\right)+C .\]
%@ b12 | div(Symbol('u')**3,Symbol('u')+1,Symbol('u')) == (Symbol('u')**2-Symbol('u')+1,-1) and ANT(2*sqrt(x)-3*x**Q(1,3)+6*x**Q(1,6)-6*log(1+x**Q(1,6)),1/(sqrt(x)+x**Q(1,3)),[Q(1,1),Q(2,1),Q(5,1),Q(64,1)]) == 0 | =2\sqrt x-3\sqrt[3]x+6\sqrt[6]x-6\ln\left(1+\sqrt[6]x\right)+C .\]
\emph{Check at $x=64$} ($u=2$): the integrand is $\frac1{8+4}=\frac1{12}$. The derivative of the answer is $\frac{du}{dx}\cdot6\left(u^2-u+1-\frac1{u+1}\right)$ with $\frac{du}{dx}=\frac1{6u^5}=\frac1{192}$ and $6\left(4-2+1-\frac13\right)=16$, giving $\frac{16}{192}=\frac1{12}$.
%@ b12b | Q(1,6*2**5) == Q(1,192) and 6*(4-2+1-Q(1,3)) == 16 and Q(16,192) == Q(1,12) | giving $\frac{16}{192}=\frac1{12}$.
\end{exmp}

\subsection*{Group 7: Proofs}

\begin{exmp}{A symmetry trick for definite integrals}{b13}
Prove that if $f$ is continuous on $[0,a]$, then $\int_0^af(x)\,dx=\int_0^af(a-x)\,dx$. Use it to evaluate $\displaystyle\int_0^{\pi/2}\frac{\sin x}{\sin x+\cos x}\,dx$ without finding an antiderivative.

\textbf{Plan.} Substitute $u=a-x$ in the right-hand integral and use Theorem~\ref{8.1-thm:subst}(2), taking care of the limits and the sign of $du$.

\textbf{Proof.} Let $u=g(x)=a-x$; $g$ has the continuous derivative $g'(x)=-1$. The function $u\mapsto f(u)$ is continuous on $[0,a]$, which contains $g([0,a])$. By Theorem~\ref{8.1-thm:subst}(2),
\[\int_0^af(a-x)\,dx=-\int_0^af(g(x))\,g'(x)\,dx=-\int_{g(0)}^{g(a)}f(u)\,du=-\int_a^0f(u)\,du=\int_0^af(u)\,du .\]
The last step reverses the limits, which changes the sign; the name of the variable does not matter. $\blacksquare$

\emph{Common mistake:} forgetting the minus sign from $g'(x)=-1$, or keeping the old limits $0$ and $a$ after the substitution (they become $a$ and $0$).

\emph{Application.} Let $I=\int_0^{\pi/2}\frac{\sin x}{\sin x+\cos x}\,dx$ and take $a=\frac\pi2$, $f(x)=\frac{\sin x}{\sin x+\cos x}$. Since $\sin\left(\frac\pi2-x\right)=\cos x$ and $\cos\left(\frac\pi2-x\right)=\sin x$, we have $f\left(\frac\pi2-x\right)=\frac{\cos x}{\cos x+\sin x}$, so $I=\int_0^{\pi/2}\frac{\cos x}{\sin x+\cos x}\,dx$. Add the two expressions for $I$: $2I=\int_0^{\pi/2}\frac{\sin x+\cos x}{\sin x+\cos x}\,dx=\int_0^{\pi/2}1\,dx=\frac\pi2$. Therefore $I=\frac\pi4$, the value found in Example~\ref{ex:m11}.
%@ b13 | simplify(sin(pi/2-x)/(sin(pi/2-x)+cos(pi/2-x))-cos(x)/(sin(x)+cos(x))) == 0 and NUM(Integral(cos(x)/(sin(x)+cos(x)),(x,0,pi/2)),pi/4) == 0 and NUM(Integral(sin(x)/(sin(x)+cos(x)),(x,0,pi/2)),pi/4) == 0 | Therefore $I=\frac\pi4$, the value found in Example~\ref{ex:m11}.
\end{exmp}

\begin{exmp}{A reduction formula for $\int(\ln x)^n\,dx$}{b14}
Prove that for $x>0$ and every integer $n\ge1$, with $I_n=\int(\ln x)^n\,dx$,
\[I_n=x(\ln x)^n-n\,I_{n-1},\qquad I_0=x+C .\]
Then find $\displaystyle\int(\ln x)^3\,dx$.

\textbf{Plan.} Parts with $p=(\ln x)^n$ and $dv=dx$, as for $\ln x$ itself in Section~8.2.

\textbf{Proof.} Let $p=(\ln x)^n$ and $dv=dx$; then $dp=n(\ln x)^{n-1}\cdot\frac1x\,dx$ (Chain Rule) and $v=x$. The parts formula gives
\[I_n=x(\ln x)^n-\int x\cdot n(\ln x)^{n-1}\frac1x\,dx=x(\ln x)^n-n\int(\ln x)^{n-1}\,dx=x(\ln x)^n-nI_{n-1}.\]
$I_0=\int1\,dx=x+C$. $\blacksquare$

\emph{Common mistake:} the factor $x$ from $v$ cancels the $\frac1x$ in $dp$; students sometimes lose the factor $n$ from the Chain Rule.

\emph{Application, $n=3$.} $I_1=x\ln x-I_0=x\ln x-x$. $I_2=x(\ln x)^2-2I_1=x(\ln x)^2-2x\ln x+2x$. $I_3=x(\ln x)^3-3I_2$, so
\[\int(\ln x)^3\,dx=x(\ln x)^3-3x(\ln x)^2+6x\ln x-6x+C .\]
%@ b14 | ANT(x*log(x)**3-3*x*log(x)**2+6*x*log(x)-6*x,log(x)**3,[Q(1,2),Q(3,2),Q(5,1),Q(10,1)]) == 0 and ANT(x*log(x)**2-2*x*log(x)+2*x,log(x)**2,[Q(1,2),Q(3,2),Q(5,1)]) == 0 | =x(\ln x)^3-3x(\ln x)^2+6x\ln x-6x+C .\]
\end{exmp}

% =====================================================================
\subsection*{Answers and discussion}

\begin{description}
  \item[Knowledge Check \ref{kc:kc1}.] (a) Substitution $u=x^3$: $x^2\,dx=\frac13\,du$, giving $\frac13e^{x^3}+C$ (no parts needed). (b) Parts with $p=x$, $dv=e^{3x}dx$: $\frac x3e^{3x}-\frac13\int e^{3x}dx=\left(\frac x3-\frac19\right)e^{3x}+C$. (c) The numerator is half the derivative of the denominator ($\frac12(2x+2)=x+1$), so $u=x^2+2x+5$ gives $\frac12\ln\left(x^2+2x+5\right)+C$; no trigonometric substitution or partial fractions. (d) The root $\sqrt{x^2+9}$ has the form $\sqrt{a^2+x^2}$ and no plain substitution is visible, so a trigonometric substitution $x=3\tan\theta$ (Example~\ref{ex:b9}).
%@ ans-kc1 | ANT(exp(x**3)/3,x**2*exp(x**3),[Q(1,2),Q(1,1),-Q(1,1)]) == 0 and ANT((x/3-Q(1,9))*exp(3*x),x*exp(3*x),[Q(1,2),Q(1,1),-Q(1,1)]) == 0 and ANT(log(x**2+2*x+5)/2,(x+1)/(x**2+2*x+5),[Q(1,2),Q(1,1),-Q(3,1)]) == 0 | so a trigonometric substitution $x=3\tan\theta$ (Example~\ref{ex:b9}).

  \item[Knowledge Check \ref{kc:kc2}.] (a) $u=\sqrt{3x+2}$, $x=\frac{u^2-2}3$, $dx=\frac{2u}3\,du$, giving $\int\frac{2u}{3(1+u)}\,du$. (b) $u=e^x$, $dx=\frac{du}u$, giving $\int\frac{u^2}{u+1}\cdot\frac1u\,du=\int\frac u{u+1}\,du$. (c) The roots have indices $2$ and $3$; use $x=u^6$, $dx=6u^5\,du$, $\sqrt x=u^3$, $\sqrt[3]x=u^2$, giving $\int\frac{u^3}{1+u^2}\cdot6u^5\,du=\int\frac{6u^8}{1+u^2}\,du$, an improper rational function.
%@ ans-kc2 | simplify(diff((Symbol('u')**2-2)/3,Symbol('u'))/(1+Symbol('u'))-2*Symbol('u')/(3*(1+Symbol('u')))) == 0 and simplify(Symbol('u')**2/(Symbol('u')+1)/Symbol('u')-Symbol('u')/(Symbol('u')+1)) == 0 and simplify(Symbol('u')**3/(1+Symbol('u')**2)*6*Symbol('u')**5-6*Symbol('u')**8/(1+Symbol('u')**2)) == 0 | giving $\int\frac{u^3}{1+u^2}\cdot6u^5\,du=\int\frac{6u^8}{1+u^2}\,du$, an improper rational function.

  \item[Knowledge Check \ref{kc:kc3}.] (a) Yes: $u=x^2$ gives $\frac12e^{x^2}+C$ (the derivative of the exponent is present). (b) No: Theorem~\ref{thm:conseq}. (c) Yes: $u=\sin x$ gives $e^{\sin x}+C$ (again the derivative of the exponent, $\cos x$, is present). (d) No: Theorem~\ref{thm:conseq}, which rests on Theorem~\ref{thm:norat}(b) and the quoted Theorem~\ref{thm:liouville}.
%@ ans-kc3 | ANT(exp(x**2)/2,x*exp(x**2),[Q(1,2),Q(1,1),-Q(3,2)]) == 0 and ANT(exp(sin(x)),cos(x)*exp(sin(x)),[Q(1,2),Q(1,1),-Q(3,2)]) == 0 | (again the derivative of the exponent, $\cos x$, is present).

  \item[Pause and think \ref{think:deadend}.] Parts helps when differentiating $u$ and integrating $dv$ make the integrand \emph{simpler}. Here neither does: the derivative of $e^{x^3}$ is $3x^2e^{x^3}$, as complicated as before, and integrating $x^5$ raises the power from $x^5$ to $x^6$, so the new integrand has $x^8$. Before integrating, look for a function and its derivative: $x^5e^{x^3}=x^3\,e^{x^3}\cdot x^2$ contains $g=x^3$ and $g'=3x^2$, so a substitution comes first, and only afterwards is the product $ue^u$ handled by parts. The rule of thumb is to try substitution before parts whenever an inside function with its derivative is visible.
%@ ans-deadend | simplify(x**5*exp(x**3)-x**3*exp(x**3)*x**2) == 0 and diff(x**3,x) == 3*x**2 | so a substitution comes first, and only afterwards is the product $ue^u$ handled by parts.

  \item[Pause and think \ref{think:twoanswers}.] With $u=\sin x$: $\int\sin x\cos x\,dx=\frac{\sin^2x}2+C$. With $u=\cos x$, $du=-\sin x\,dx$: $-\int u\,du=-\frac{\cos^2x}2+C$. Both are antiderivatives of $\sin x\cos x$ on all of $\mathbb R$, and they are not equal as functions; their difference is $\frac{\sin^2x}2-\left(-\frac{\cos^2x}2\right)=\frac{\sin^2x+\cos^2x}2=\frac12$, a constant, exactly as Theorem~\ref{thm:check}(2) says. (A third form, $-\frac14\cos2x$, differs from the first by $\frac14$.)
%@ ans-twoanswers | ANT(sin(x)**2/2,sin(x)*cos(x),[Q(1,2),Q(1,1),Q(3,1)]) == 0 and ANT(-cos(x)**2/2,sin(x)*cos(x),[Q(1,2),Q(1,1),Q(3,1)]) == 0 and ANT(-cos(2*x)/4,sin(x)*cos(x),[Q(1,2),Q(1,1),Q(3,1)]) == 0 and simplify(sin(x)**2/2-(-cos(2*x)/4)-Q(1,4)) == 0 and simplify(sin(x)**2/2+cos(x)**2/2) == Q(1,2) | differs from the first by $\frac14$.)

  \item[Pause and think \ref{think:provenot}.] A failed search proves only that those methods failed; a cleverer substitution might exist, and mathematicians did not know for a long time. What proves the claim is a theorem: Liouville's criterion says that if $e^{-x^2}$ had an elementary antiderivative, there would be a rational $R$ with $R'-2xR=1$, and Theorem~\ref{thm:norat}(a) shows by an algebraic argument (using only Euclid's lemma and degrees) that no such $R$ exists. We proved the algebra in full; we did \emph{not} prove Liouville's criterion (Theorem~\ref{thm:liouville}), which is quoted.
\end{description}

\end{document}
